% Generated draft: requires mathematical and source review.
\documentclass[11pt]{article}
\usepackage[margin=25mm]{geometry}
\usepackage{fontspec}
\setmainfont{Latin Modern Roman}
\usepackage{amsmath,amssymb,mathtools,unicode-math}
\setmathfont{Latin Modern Math}
\usepackage{enumitem,needspace,xurl,hyperref}
\setlength{\parindent}{0pt}
\setlength{\parskip}{5pt}
\setlength{\emergencystretch}{3em}
\newcommand{\N}{\mathbb{N}}
\newcommand{\Z}{\mathbb{Z}}
\newcommand{\Q}{\mathbb{Q}}
\newcommand{\R}{\mathbb{R}}
\newcommand{\C}{\mathbb{C}}
\newcommand{\F}{\mathbb{F}}
\newcommand{\A}{\mathbb{A}}
\newcommand{\PP}{\mathbb{P}}
\newcommand{\E}{\mathbb{E}}
\newcommand{\eps}{\varepsilon}
\newcommand{\dd}{\,\mathrm{d}}
\begin{document}
\section*{UnsolvedMath ID: 30001609}
\textbf{Source problem code:} catalogue entry
\subsection{Definitions and mathematical statement}
The supplied question concerns whether every model space admits an unconditional basis of reproducing kernels. The Hardy-space domain and the permitted locations of the kernels are not specified in the record. The following is a precise formulation under the explicit convention of the unit disc and interior evaluation points; this convention requires source confirmation.

Let \(\mathbb D=\{z\in\C:|z|<1\}\), and define
\[
H^2(\mathbb D)=\left\{f(z)=\sum_{n=0}^{\infty}a_nz^n:\sum_{n=0}^{\infty}|a_n|^2<\infty\right\},
\qquad \|f\|_{H^2}^2=\sum_{n=0}^{\infty}|a_n|^2.
\]
An inner function \(\Theta\) is a bounded analytic function on \(\mathbb D\) whose radial boundary values have modulus one almost everywhere on the unit circle with respect to arc length. Its model space is
\[
K_\Theta=H^2(\mathbb D)\ominus\Theta H^2(\mathbb D).
\]
For \(\lambda\in\mathbb D\), its reproducing kernel is
\[
k_\lambda^\Theta(z)=\frac{1-\overline{\Theta(\lambda)}\Theta(z)}{1-\overline\lambda z},
\qquad f(\lambda)=\langle f,k_\lambda^\Theta\rangle_{H^2}\quad(f\in K_\Theta),
\]
where the inner product is linear in its first argument. For nonconstant \(\Theta\), put
\[
e_\lambda^\Theta=\frac{k_\lambda^\Theta}{\|k_\lambda^\Theta\|},
\qquad \|k_\lambda^\Theta\|^2=\frac{1-|\Theta(\lambda)|^2}{1-|\lambda|^2}.
\]
An unconditional basis is a family such that every vector has a unique expansion in that family, with convergence in norm independent of the order of summation.

\textbf{Question under the stated convention.} Is it true that for every nonconstant inner function \(\Theta\) there exist a finite or countably infinite index set \(J\) and distinct points \(\lambda_j\in\mathbb D\), \(j\in J\), such that \((e_{\lambda_j}^\Theta)_{j\in J}\) is an unconditional basis of \(K_\Theta\)? Equivalently, is there such a family with dense linear span in \(K_\Theta\) and constants \(0<A\leq B<\infty\) satisfying
\[
A\sum_{j\in J}|c_j|^2
\leq \left\|\sum_{j\in J}c_j e_{\lambda_j}^\Theta\right\|_{H^2}^2
\leq B\sum_{j\in J}|c_j|^2
\]
for every finitely supported complex family \((c_j)_{j\in J}\)? The points and constants may depend on \(\Theta\). Normalization does not change the unconditional-basis question for nonzero kernels. For a constant inner function, \(K_\Theta=\{0\}\), which may be handled by the empty-basis convention.

The source does not establish whether kernels at boundary points admitting bounded evaluation are also allowed, or whether another Hardy-space domain is intended. These choices must be checked before this conditional formulation is accepted as the intended statement.
\subsection{Short English statement}
The supplied problem asks whether every model space \(H^2\ominus\Theta H^2\) has an unconditional basis consisting of reproducing kernels. The displayed formulation uses the unit disc and interior kernels explicitly; the source record leaves the domain and admissible kernel locations unresolved. Its open status is supplied dated curation, not an independently verified conclusion.
\subsection{Sources}
\begin{itemize}
\item Supplied local-registry record, permanent UnsolvedMath ID \(30001609\), source problem code \texttt{catalogue entry}, title: \emph{Unconditional Reproducing-Kernel Bases in Model Spaces}. Project registry references: \texttt{data/unsolvedmath/problems.json} and \texttt{catalog/problems.json}; no independent file inspection was performed.
\item Supplied source citation: \emph{Operator Theory and Harmonic Analysis (2011)}, Oberwolfach Reports, pp.~2813--2875, DOI \texttt{10.4171/owr/2010/49}. Supplied source URL and literature-source link labelled ``original problem report'': \url{https://doi.org/10.4171/owr/2010/49}. No individual proposer is supplied. The record gives the proposed year as 2010.
\item Literature assessment dated exactly \texttt{2026-08-21}: the supplied curation labels the problem open, says that the report explicitly presents the general existence question as open, and reports that later work provides criteria and examples without a general resolution being found. This is supplied dated curation, not independently verified literature assessment; the background explicitly requests independent verification of time-sensitive status and cited claims.
\item The supplied statement status is \texttt{corrected\_verified}, accompanied by a report of manual semantic, typographic, and LaTeX review. This supplied editorial label does not independently establish the intended domain, boundary-kernel convention, or mathematical status.
\end{itemize}
\end{document}
